\textbf{Время:} 01:30

\textbf{Тезис:} при коллективном расширении пары с выбранным импульсом характеризуют область однородности, а не весь источник.

\textbf{На экране:} сначала весь светлый источник и расходящиеся стрелки; затем выделенный золотой участок и стрелки K\_T вправо; после этого две фразы справа, завершая термином «область однородности».

\textbf{Текст:} На предыдущем рисунке мы выбирали пары из определённого интервала среднего импульса. Теперь посмотрим, почему этот отбор важен для пространственного размера.

Светлая область слева — весь расширяющийся источник. Расходящиеся стрелки показывают среднее движение вещества. В разных частях системы оно направлено по-разному. Поэтому средний импульс испущенных частиц связан с местом испускания — это первая фраза справа.

Выберем пары, летящие вправо. На рисунке их направление отмечено золотыми стрелками. Такие пары преимущественно приходят из участка, где движение вещества согласовано с выбранным направлением и величиной импульса. Этот участок выделен золотым. Его называют областью однородности.

Это не жёстко отрезанный кусок системы: тепловое движение размывает границы и позволяет частицам приходить из разных мест. Но основная чувствительность корреляции приходится на область, выбранную импульсами пары.

При увеличении k-T эта область может становиться меньше, хотя весь источник не сжимается. Поэтому зависимость радиусов от импульса помогает исследовать коллективное расширение.

\textbf{Переход:} Теперь прочитаем именно такую зависимость на следующем графике STAR.

\textbf{Справка:} Схема качественная; она иллюстрирует корреляции координаты и импульса и не устанавливает единственную динамическую модель. \(k_T=|\mathbf K_T|\) — модуль среднего поперечного импульса пары. HBT-радиусы характеризуют области однородности при выбранных K и условиях регистрации; они не тождественны полному геометрическому радиусу системы. Уменьшение радиусов с k\_T совместимо с коллективным потоком, но не служит его единственным возможным объяснением без дополнительных наблюдаемых. Lisa et al. (2005), §2.7; Краева (2024), §3–4.
